% Generated draft: requires mathematical and source review.
\documentclass[11pt]{article}
\usepackage[margin=25mm]{geometry}
\usepackage{fontspec}
\setmainfont{Latin Modern Roman}
\usepackage{amsmath,amssymb,mathtools,unicode-math}
\setmathfont{Latin Modern Math}
\usepackage{enumitem,needspace,xurl,hyperref}
\setlength{\parindent}{0pt}
\setlength{\parskip}{5pt}
\setlength{\emergencystretch}{3em}
\newcommand{\N}{\mathbb{N}}
\newcommand{\Z}{\mathbb{Z}}
\newcommand{\Q}{\mathbb{Q}}
\newcommand{\R}{\mathbb{R}}
\newcommand{\C}{\mathbb{C}}
\newcommand{\F}{\mathbb{F}}
\newcommand{\A}{\mathbb{A}}
\newcommand{\PP}{\mathbb{P}}
\newcommand{\E}{\mathbb{E}}
\newcommand{\eps}{\varepsilon}
\newcommand{\dd}{\,\mathrm{d}}
\begin{document}
\section*{UnsolvedMath ID: 30006643}
\textbf{Source problem code:} catalogue entry
\subsection{Definitions and mathematical statement}
Let \(\pi\) denote the non-Gaussian log-concave target in the supplied record, and let \(\lambda>0\) denote its stated convexity parameter. Write \(G\) for the source's admissible parameter space, \(E\colon G\to\R\) for its energy, and \((\Sigma(t),m(t))\) for its projected gradient trajectory with initial state \((\Sigma_0,m_0)\). These symbols identify the objects mentioned in the excerpt; the target's underlying space, the precise convexity hypothesis, the definitions of \(G\) and \(E\), and the evolution equation are not supplied.

As background, the excerpt reports Theorem 5 (Unique minimizer): \(E\) attains a unique minimizer \((\Sigma_{\mathrm{opt}},m_{\mathrm{opt}})\in G\). Set
\[
E_*:=\min_{(\Lambda,n)\in G}E(\Lambda,n)
     =E(\Sigma_{\mathrm{opt}},m_{\mathrm{opt}}).
\]
The excerpt also reports Theorem 6 (Convergence), whose displayed estimate can be written equivalently as
\[
E(\Sigma(t),m(t))-E_*
\leq e^{-\gamma t}\bigl(E(\Sigma_0,m_0)-E_*\bigr),
\qquad \gamma\propto\lambda.
\]
This is a reported background result, subject to hypotheses absent from the supplied excerpt. In particular, the admissible initial states, time interval, and dependencies of the proportionality constants are unspecified.

The requested problem has two parts:
\begin{enumerate}
\item For the source's projected gradient dynamics with a non-Gaussian log-concave target, establish a sharper decay estimate whose convergence guarantee does not deteriorate with \(\lambda\). The excerpt motivates this request by comparison with the Gaussian-target case. It does not specify a replacement rate or the quantities to be held fixed when varying \(\lambda\), so a precise uniform inequality cannot be determined from this record alone.
\item Construct a numerical scheme for these dynamics and prove its convergence for non-Gaussian targets under explicitly stated assumptions. The source does not specify the discretization, convergence topology, step-size conditions, or whether convergence is intended toward the continuous trajectory, the unique energy minimizer, or both.
\end{enumerate}
Thus this is a statement-only formulation of the supplied research objectives; the missing model definitions and theorem hypotheses must be recovered before a fully quantified mathematical statement can be given.
\subsection{Short English statement}
For the projected gradient model described by the source, improve the energy decay guarantee for non-Gaussian log-concave targets so that it does not deteriorate with the convexity parameter \(\lambda>0\), and construct a numerical method with a convergence proof. The supplied excerpt omits essential model definitions and convergence specifications, so these objectives require further source clarification.
\subsection{Sources}
\begin{itemize}
\item Supplied local UnsolvedMath registry record, permanent ID \(30006643\), title \emph{Sharper Decay and Algorithms for Non-Gaussian Targets}; exact source problem code \texttt{catalogue entry}. The record supplies no named proposer and gives proposed year 2026.
\item Originating report: \emph{Flows on Measure Spaces and Applications in Machine Learning} (2026), Oberwolfach Reports, pp.\ 795--884, DOI \texttt{10.4171/owr/2026/13}. Supplied source and literature-source link: \url{https://doi.org/10.4171/owr/2026/13}. The supplied excerpt attributes the minimizer and decay statements to Theorems 5 and 6, respectively, without their complete hypotheses or a pinpoint page.
\item Supplied literature source: \emph{Provably convergent stochastic fixed-point algorithm for Wasserstein barycenters}, \url{https://arxiv.org/abs/2505.24384}. Its relevance and applicability to the stated model have not been independently verified here.
\item Literature assessment dated exactly \texttt{2026-08-22}: the supplied curation labels the problem open and states that the report leaves both the sharper non-Gaussian decay estimate and a good provably convergent numerical scheme open, while available algorithms for related Wasserstein problems do not prove these claims for the stated model. This is supplied dated curation, not an independently verified literature assessment; the record itself requests independent verification of time-sensitive status and cited claims.
\item The registry labels the statement \texttt{corrected\_verified} and says it was reconstructed and checked against the official Oberwolfach report. This is supplied verification metadata, not a source check performed for this draft.
\end{itemize}
\end{document}
